\section{AMC 10 10B 2010}
        \begin{enumerate}
        	\item (\textit{AMC 10 10B 2010}) What is $100(100-3)-(100\cdot100-3)$?


 $\textbf{(A)}\ -20,000 \qquad \textbf{(B)}\ -10,000 \qquad \textbf{(C)}\ -297 \qquad \textbf{(D)}\ -6 \qquad \textbf{(E)}\ 0$


 
 \textbf{Answer: }C



\textbf{Solution}

$100(100-3)-(100\cdot{100}-3)=10000-300-10000+3=-300+3=\boxed{\textbf{(C)}\ -297}$.

 


	\item (\textit{AMC 10 10B 2010}) Makayla attended two meetings during her $9$-hour work day. The first meeting took $45$ minutes and the second meeting took twice as long. What percent of her work day was spent attending meetings?


 $\textbf{(A)}\ 15 \qquad \textbf{(B)}\ 20 \qquad \textbf{(C)}\ 25 \qquad \textbf{(D)}\ 30 \qquad \textbf{(E)}\ 35$


 
 \textbf{Answer: }C



\textbf{Solution}

The total number of minutes in her $9$-hour work day is $9 \times 60 = 540.$The total amount of time spend in meetings in minutes is $45 + 45 \times 2 = 135.$The answer is then $\frac{135}{540}$ $= \boxed{25\%}$ or  $\boxed{(C)}$

 


	\item (\textit{AMC 10 10B 2010}) A drawer contains red, green, blue, and white socks with at least 2 of each color. What is
the minimum number of socks that must be pulled from the drawer to guarantee a matching
pair?


 $\textbf{(A)}\ 3 \qquad \textbf{(B)}\ 4 \qquad \textbf{(C)}\ 5 \qquad \textbf{(D)}\ 8 \qquad \textbf{(E)}\ 9$


 
 \textbf{Answer: }C



\textbf{Solution}

After you draw $4$ socks, you can have one of each color, so (according to the \href{https://artofproblemsolving.com/wiki/index.php?title=Pigeonhole\_principle}{pigeonhole principle}), if you pull $\boxed{\textbf{(C)}\ 5}$ then you will be guaranteed a matching pair.

 


	\item (\textit{AMC 10 10B 2010}) For a real number $x$, define $\heartsuit(x)$ to be the average of $x$ and $x^2$. What is $\heartsuit(1)+\heartsuit(2)+\heartsuit(3)$?


 $\textbf{(A)}\ 3 \qquad \textbf{(B)}\ 6 \qquad \textbf{(C)}\ 10 \qquad \textbf{(D)}\ 12 \qquad \textbf{(E)}\ 20$


 
 \textbf{Answer: }C



\textbf{Solution}

The average of two numbers, $a$ and $b$, is defined as $\frac{a+b}{2}$. Thus the average of $x$ and $x^2$ would be $\frac{x(x+1)}{2}$. With that said,  we need to find the sum when we plug, $1$, $2$ and $3$ into that equation. So:

 

 \[\frac{1(1+1)}{2} + \frac{2(2+1)}{2} + \frac{3(3+1)}{2} = \frac{2}{2} + \frac{6}{2} + \frac{12}{2} = 1+3+6= \boxed{\textbf{(C)} 10}\]

 


	\item (\textit{AMC 10 10B 2010}) A month with $31$ days has the same number of Mondays and Wednesdays. How many of the seven days of the week could be the first day of this month?


 $\textbf{(A)}\ 2 \qquad \textbf{(B)}\ 3 \qquad \textbf{(C)}\ 4 \qquad \textbf{(D)}\ 5 \qquad \textbf{(E)}\ 6$


 
 \textbf{Answer: }B



\textbf{Solution}

$31 \equiv 3 \pmod {7}$ so the week cannot start with Saturday, Sunday, Tuesday or Wednesday as that would result in an unequal number of Mondays and Wednesdays. Therefore, Monday, Thursday, and Friday are valid so the answer is $\boxed{B}$.

 


	\item (\textit{AMC 10 10B 2010}) A circle is centered at $O$, $\overline{AB}$ is a diameter and $C$ is a point on the circle with $\angle COB = 50^\circ$.
What is the degree measure of $\angle CAB$?


 $\textbf{(A)}\ 20 \qquad \textbf{(B)}\ 25 \qquad \textbf{(C)}\ 45 \qquad \textbf{(D)}\ 50 \qquad \textbf{(E)}\ 65$


 
 \textbf{Answer: }B



\textbf{Solution 1}

Assuming we do not already know an inscribed angle is always half of its central angle, we will try a different approach. Since $O$ is the center, $OC$ and $OA$ are radii and they are congruent. Thus, $\triangle COA$ is an isosceles triangle. Also, note that $\angle COB$ and $\angle COA$ are supplementary, then $\angle COA = 180 - 50 = 130^{\circ}$. Since $\triangle COA$ is isosceles, then $\angle OCA \cong \angle OAC$. They also sum to $50^{\circ}$, so each angle is $\boxed{\textbf{(B)}\ 25}$.

 


 



\textbf{Solution 2 (Alcumus)}

Note that $\angle AOC = 180^\circ - 50^\circ = 130^\circ$. Because triangle $AOC$ is isosceles, $\angle CAB = (180^\circ - 130^\circ)/2 = \boxed{25^\circ}$.

 
                    \begin{figure}[H]
                        \centering
                        \includegraphics[width=0.5\textwidth]{Solution/AMC_10_10B_2010_6_Solution_Figure_1.png}
                    \end{figure}
                    

 


	\item (\textit{AMC 10 10B 2010}) A triangle has side lengths $10$, $10$, and $12$. A rectangle has width $4$ and area equal to the
area of the triangle. What is the perimeter of this rectangle?


 $\textbf{(A)}\ 16 \qquad \textbf{(B)}\ 24 \qquad \textbf{(C)}\ 28 \qquad \textbf{(D)}\ 32 \qquad \textbf{(E)}\ 36$


 
 \textbf{Answer: }D



\textbf{Solution 1}

The triangle is isosceles. The height of the triangle is therefore given by $h = \sqrt{10^2 - \left( \dfrac{12}{2} \right)^2} =  \sqrt{64} = 8$

 Now, the area of the triangle is $\dfrac{bh}{2} = \dfrac{12*8}{2} = \dfrac{96}{2} = 48$

 We have that the area of the rectangle is the same as the area of the triangle, namely $48$. We also have the width of the rectangle: $4$.

 The length of the rectangle therefore is:$l = \dfrac{48}{4} = 12$

 The perimeter of the rectangle then becomes:$2l + 2w = 2*12 + 2*4 = 32$

 The answer is:

 $\boxed{\textbf{(D)}\ 32}$

 



\textbf{Solution 2}

An alternative way to find the area of the triangle is by using Heron's formula, $A=\sqrt{(s)(s-a)(s-b)(s-c)}$ where $s$ is the semi-perimeter of the triangle (meaning half the perimeter). Here, the semi-perimeter is $(10+10+12)/2 = 16$. Thus the area equals $\sqrt{(16)(16-10)(16-10)(16-12)} = \sqrt{16*6*6*4} = 48.$ We know that the width of the rectangle is $4$, so $48/4 = 12$, which is the length. The perimeter of the rectangle is $2(4+12) =$ $\boxed{\textbf{(D)}\ 32}$.

 (Solution by Flamedragon)

 



\textbf{Solution 3}

Note that a triangle with side lengths $10,10,12$ is essentially $2$   “6,8,10” right triangles stuck together. Hence, the height is $8$, and our area is $48$.

 So, the length of the rectangle is $\frac{48}{4}=12$, and our perimeter $P=2(4+12)=\boxed{\textbf{(D)}\ 32}$

 


	\item (\textit{AMC 10 10B 2010}) A ticket to a school play cost $x$ dollars, where $x$ is a whole number. A group of 9th graders buys tickets costing a total of $\textdollar 48$, and a group of 10th graders buys tickets costing a total of $\textdollar 64$. How many values for $x$ are possible?


 $\textbf{(A)}\ 1 \qquad \textbf{(B)}\ 2 \qquad \textbf{(C)}\ 3 \qquad \textbf{(D)}\ 4 \qquad \textbf{(E)}\ 5$


 
 \textbf{Answer: }E



\textbf{Solution}

We find the greatest common factor of $48$ and $64$ to be $16$. The number of factors of $16$ is $5$ which is the answer $(E)$.

 



\textbf{Solution 2}

Since the price of one ticket is a whole number, it must be a divisor of both $48$ and $64$. An integer is a common divisor of two numbers if and only if it is divisible by their difference. Hence, all the possible prices are positive divisors of $16$, or $1$, $2$, $4$, $8$, and $16$, making the answer $5$, or E. - By Monkey\_King

 


	\item (\textit{AMC 10 10B 2010}) Lucky Larry's teacher asked him to substitute numbers for $a$, $b$, $c$, $d$, and $e$ in the expression $a-(b-(c-(d+e)))$ and evaluate the result. Larry ignored the parentheses but added and subtracted correctly and obtained the correct result by coincidence. The numbers Larry substituted for $a$, $b$, $c$, and $d$ were $1$, $2$, $3$, and $4$, respectively. What number did Larry substitute for $e$?


 $\textbf{(A)}\ -5 \qquad \textbf{(B)}\ -3 \qquad \textbf{(C)}\ 0 \qquad \textbf{(D)}\ 3 \qquad \textbf{(E)}\ 5$


 
 \textbf{Answer: }D



\textbf{Solution}

We simply plug in the numbers
 \[1 - 2 - 3 - 4 + e = 1 - (2 - (3 - (4 + e)))\]
 \[-8 + e = -2 - e\]
 \[2e = 6\]
 \[e = 3 \;\;(D)\]

 


	\item (\textit{AMC 10 10B 2010}) Shelby drives her scooter at a speed of $30$ miles per hour if it is not raining, and $20$ miles per hour if it is raining. Today she drove in the sun in the morning and in the rain in the evening, for a total of $16$ miles in $40$ minutes. How many minutes did she drive in the rain?


 $\textbf{(A)}\ 18 \qquad \textbf{(B)}\ 21 \qquad \textbf{(C)}\ 24 \qquad \textbf{(D)}\ 27 \qquad \textbf{(E)}\ 30$


 
 \textbf{Answer: }C



\textbf{Solution 1}

Let $x$ be the time it is not raining, and $y$ be the time it is raining, in hours.

 We have the system: $30x+20y=16$ and $x+y=\frac{2}{3}$

 Solving gives $x=\frac{4}{15}$ and $y=\frac{2}{5}$

 We want $y$ in minutes, $\frac{2}{5} \cdot 60=24 \Rightarrow C$

 



\textbf{Solution 2}

Let $x$ be the time she drove in the rain. Thus, the number of minutes she did not drive in the rain is $40-x$ . 

 Since we are calculating the time in minutes, it is best to convert the speeds in minutes. Thus, the speed per minute when it is not raining is $\frac{1}{2}$ miles per minute, and $\frac{1}{3}$ miles per minute when it is raining. Thus, we have the equation, $\frac{1}{2} \cdot (40-x) + \frac{1}{3} \cdot x = 16$ 

 Solving, gives $x$ = $24$ , so the amount of time she drove in the rain is $24$ minutes.

 


	\item (\textit{AMC 10 10B 2010}) A shopper plans to purchase an item that has a listed price greater than $\textdollar 100$ and can use any one of the three coupons. Coupon A gives $15\%$ off the listed price, Coupon B gives $\textdollar 30$ off the listed price, and Coupon C gives $25\%$ off the amount by which the listed price exceeds
$\textdollar 100$. 

Let $x$ and $y$ be the smallest and largest prices, respectively, for which Coupon A saves at least as many dollars as Coupon B or Coupon C. What is $y - x$?


 $\textbf{(A)}\ 50 \qquad \textbf{(B)}\ 60 \qquad \textbf{(C)}\ 75 \qquad \textbf{(D)}\ 80  \qquad \textbf{(E)}\ 100$


 
 \textbf{Answer: }A



\textbf{Solution}

Let the listed price be $(100 + p)$, where $p > 0$

 Coupon A saves us: $0.15(100+p) = (0.15p + 15)$

 Coupon B saves us: $30$

 Coupon C saves us: $0.25p$

 Now, the condition is that A has to be greater than or equal to either B or C which gives us the following inequalities:

 $A \geq B \Rightarrow 0.15p + 15 \geq 30 \Rightarrow p \geq 100$

 $A \geq C \Rightarrow 0.15p + 15 \geq 0.25p \Rightarrow p \leq 150$

 We see here that the greatest possible value for $p$ is $150$, thus $y = 100 + 150 = 250$ and the smallest value for $p$ is $100$ so $x = 100 + 100 = 200$. 

 The difference between $y$ and $x$ is $y - x = 250 - 200 = \boxed{\textbf{(A)}\ 50}$

 


	\item (\textit{AMC 10 10B 2010}) At the beginning of the school year, $50\%$ of all students in Mr. Wells' math class answered "Yes" to the question "Do you love math", and $50\%$ answered "No." At the end of the school year, $70\%$ answered "Yes" and $30\%$ answered "No." Altogether, $x\%$ of the students gave a different answer at the beginning and end of the school year. What is the difference between the maximum and the minimum possible values of $x$?


 $\textbf{(A)}\ 0 \qquad \textbf{(B)}\ 20 \qquad \textbf{(C)}\ 40 \qquad \textbf{(D)}\ 60 \qquad \textbf{(E)}\ 80$


 
 \textbf{Answer: }D



\textbf{Solution}

The minimum possible value would be $70 - 50 = 20\%$. The maximum possible value would be $30 + 50 = 80\%$. The difference is $80 - 20 = \boxed{\textbf{(D) }60}$.

 


	\item (\textit{AMC 10 10B 2010}) What is the sum of all the solutions of $x = \left|2x-|60-2x|\right|$?


 $\textbf{(A)}\ 32 \qquad \textbf{(B)}\ 60 \qquad \textbf{(C)}\ 92 \qquad \textbf{(D)}\ 120 \qquad \textbf{(E)}\ 124$


 
 \textbf{Answer: }C



\textbf{Solution 1}

We evaluate this in cases:

 \textit{Case 1}$x<30$ 

 When $x<30$ we are going to have $60-2x>0$.  When $x>0$ we are going to have $|x|>0\implies x>0$ and when $-x>0$ we are going to have $|x|>0\implies -x>0$.  Therefore we have $x=|2x-(60-2x)|$.$x=|2x-60+2x|\implies x=|4x-60|$

 \textit{Subcase 1 }$30>x>15$

 When $30>x>15$ we are going to have $4x-60>0$. When this happens, we can express $|4x-60|$ as $4x-60$.Therefore we get $x=4x-60\implies -3x=-60\implies x=20$.  We check if $x=20$ is in the domain of the numbers that we put into this subcase, and it is, since $30>20>15$. Therefore $20$ is one possible solution.  

 \textit{ Subcase 2 } $x<15$

 When $x<15$ we are going to have $4x-60<0$, therefore $|4x-60|$ can be expressed in the form $60-4x$.We have the equation $x=60-4x\implies 5x=60\implies x=12$.  Since $12$ is less than $15$, $12$ is another possible solution.$x=|2x-|60-2x||$

 \textit{Case 2 }:  $x>30$

 When $x>30$, $60-2x<0$. When $x<0$ we can express this in the form $-x$.  Therefore we have $-(60-2x)=2x-60$.  This makes sure that this is positive, since we just took the negative of a negative to get a positive.  Therefore we have $x=|2x-(2x-60)|$ 

 $x=|2x-2x+60|$

 $x=|60|$

 $x=60$

 We have now evaluated all the cases, and found the solution to be $\{60,12,20\}$ which have a sum of $\boxed{\textbf{(C)}\ 92}$

 



\textbf{Solution 2}

From the equation $x = \left|2x-|60-2x|\right|$ , we have $x = 2x-|60-2x|$ , or $-x = 2x-|60-2x|$. Therefore, $x=|60-2x|$ , or $3x=|60-2x|$. From here we have four possible cases:

 1. $x=60-2x$; this simplifies to $3x=60$, so $x=20$.

 2. $-x=60-2x$; this simplifies to $x=60$.

 3. $3x=60-2x$; this simplifies to $5x=60$, so $x=12$.

 4. $-3x=60-2x$; this simplifies to $-x=60$, so $x=-60$. However, this solution is extraneous because the absolute value of $2x-|60-2x|$ cannot be negative.

 The sum of all of the solutions of $x$ is $20+60+12=\boxed{\textbf{(C)}\ 92}$

 


	\item (\textit{AMC 10 10B 2010}) The average of the numbers $1, 2, 3,\cdots, 98, 99,$ and $x$ is $100x$. What is $x$?


 $\textbf{(A)}\ \dfrac{49}{101} \qquad \textbf{(B)}\ \dfrac{50}{101} \qquad \textbf{(C)}\ \dfrac{1}{2} \qquad \textbf{(D)}\ \dfrac{51}{101} \qquad \textbf{(E)}\ \dfrac{50}{99}$


 
 \textbf{Answer: }B



\textbf{Solution}

We first sum the first $99$ numbers: $\frac{99(100)}{2}=99\cdot50$. Then, we know that the sum of the series is $99\cdot50+x$. There are $100$ terms, so we can divide this sum by $100$ and set it equal to $100x$:
 \[\frac{99\cdot50+x}{100}=100x \Rightarrow 99\cdot50=100^2x-x\Rightarrow \frac{99\cdot50}{100^2-1}=x\]Using difference of squares:
 \[x=\frac{99\cdot50}{101\cdot99}=\frac{50}{101}\]Thus, the answer is $\boxed{\textbf{(B) }\frac{50}{101}}$.

 


	\item (\textit{AMC 10 10B 2010}) On a $50$-question multiple choice math contest, students receive $4$ points for a correct answer, $0$ points for an answer left blank, and $-1$ point for an incorrect answer. Jesse’s total score on the contest was $99$. What is the maximum number of questions that Jesse could have answered correctly?


 $\textbf{(A)}\ 25 \qquad \textbf{(B)}\ 27 \qquad \textbf{(C)}\ 29 \qquad \textbf{(D)}\ 31 \qquad \textbf{(E)}\ 33$


 
 \textbf{Answer: }C



\textbf{Solution}

Let $a$ be the amount of questions Jesse answered correctly, $b$ be the amount of questions Jesse left blank, and $c$ be the amount of questions Jesse answered incorrectly. Since there were $50$ questions on the contest, $a+b+c=50$. Since his total score was $99$, $4a-c=99$. Also, $a+c\leq50 \Rightarrow c\leq50-a$. We can substitute this inequality into the previous equation to obtain another inequality: $4a-(50-a)\leq99 \Rightarrow 5a\leq149 \Rightarrow a\leq \frac{149}5=29.8$. Since $a$ is an integer, the maximum value for $a$ is $\boxed{\textbf{(C)}\ 29}$.

 



\textbf{Solution 2}

We can plug in each answer choice and find our what is the greatest one that works. We start of with Answer Choice $E$ since it is the largest one.

 E is right, Jesse got $33$ questions right.

 If Jesse got $33$ questions right, then he gains $132$ points, he then needs to get another $33$ wrong to achieve a score of $99$. However, this is impossible as the test only contains $50$ questions, and he needs $66$ questions in order to achieve this.

 D is right, Jesse got $31$ questions right.

 If Jesse got $31$ questions right, then he gains $124$ points, he then needs to get another $25$ wrong in order to achieve a score of $99$. However, this is impossible as the test only contains $50$ questions, and he needs $56$ questions in order to achieve this.

 C is right, Jesse got $29$ questions right.

 If Jesse got $29$ questions right, then he gains $116$ points, he then needs to get another $17$ wrong in order to achieve a score of $99$. This is possible since this only requires $46$ questions. The other $4$ questions remain blank and earn him $0$ points.

 - SAMANTAP

 


	\item (\textit{AMC 10 10B 2010}) A square of side length $1$ and a circle of radius $\dfrac{\sqrt{3}}{3}$ share the same center. What is the area inside the circle, but outside the square?


 $\textbf{(A)}\ \dfrac{\pi}{3}-1 \qquad \textbf{(B)}\ \dfrac{2\pi}{9}-\dfrac{\sqrt{3}}{3} \qquad \textbf{(C)}\ \dfrac{\pi}{18} \qquad \textbf{(D)}\ \dfrac{1}{4} \qquad \textbf{(E)}\ \dfrac{2\pi}{9}$


 
 \textbf{Answer: }B



\textbf{Solution 1}

The radius of the circle is $\frac{\sqrt{3}}{3} = \sqrt{\frac{1}{3}}$. Half the diagonal of the square is $\frac{\sqrt{1^2+1^2}}{2} = \frac{\sqrt{2}}{2} = \sqrt{\frac12}$. We can see that the circle passes outside the square, but the square is NOT completely contained in the circle. Therefore the picture will look something like this:

 
                    \begin{figure}[H]
                        \centering
                        \includegraphics[width=0.5\textwidth]{Solution/AMC_10_10B_2010_16_Solution_Figure_1.png}
                    \end{figure}
                    

 Then we proceed to find: 4 $\cdot$ (area of sector marked off by the two radii - area of the triangle with sides on the square and the two radii).

 First we realize that the radius perpendicular to the side of the square between the two radii marking off the sector splits $AB$ in half. Let this half-length be $a$. Also note that $OX=\frac12$ because it is half the side length of the square. Because this is a right triangle, we can use the \href{https://artofproblemsolving.com/wiki/index.php?title=Pythagorean\_Theorem}{Pythagorean Theorem} to solve for $a.$

 
 \[a^2+\left( \frac12 \right) ^2 = \left( \frac{\sqrt{3}}{3} \right) ^2\]

 Solving, $a= \frac{\sqrt{3}}{6}$ and $2a=\frac{\sqrt{3}}{3}$. Since $AB=AO=BO$, $\triangle AOB$ is an equilateral triangle and the central angle is $60^{\circ}$. Therefore the sector has an area $\pi \left( \frac{\sqrt{3}}{3} \right) ^2 \left( \frac{60}{360} \right) = \frac{\pi}{18}$.

 Now we turn to the triangle. Since it is equilateral, we can use the formula for the \href{https://artofproblemsolving.com/wiki/index.php?title=Area\_of\_an\_equilateral\_triangle}{area of an equilateral triangle} which is

 
 \[\frac{s^2\sqrt{3}}{4} = \frac{\frac13 \sqrt{3}}{4} = \frac{\sqrt{3}}{12}\]

 Putting it together, we get the answer to be $4\cdot\left( \frac{\pi}{18}-\frac{\sqrt{3}}{12} \right)= \boxed{\textbf{(B)}\ \frac{2\pi}{9}-\frac{\sqrt{3}}{3}}$

 



\textbf{Solution 2 (Answer Choices)}

At once, we can eliminate $(C)$, $(D)$, and $(E)$ because the answer should be $a \times \pi + b$, where $a$ and $b$ are real numbers (not necessarily rational) not equal to $0$. Now, all that's left is $\frac{\pi}{3}-1$ and $\frac{2 \pi}{9} - \frac{\sqrt{3}}{3}$. $\frac{\pi}{3}-1$ is simply the area of the circle minus the square, which is definitely wrong, so the answer is $\boxed{\textbf{(B)} \frac{2 \pi}{9} - \frac{\sqrt{3}}{3}}$

 


	\item (\textit{AMC 10 10B 2010}) Every high school in the city of Euclid sent a team of $3$ students to a math contest. Each participant in the contest received a different score. Andrea's score was the median among all students, and hers was the highest score on her team. Andrea's teammates Beth and Carla placed $37$th and $64$th, respectively. How many schools are in the city?


 $\textbf{(A)}\ 22 \qquad \textbf{(B)}\ 23 \qquad \textbf{(C)}\ 24 \qquad \textbf{(D)}\ 25 \qquad \textbf{(E)}\ 26$


 
 \textbf{Answer: }B



\textbf{Solution 1}

There are $x$ schools. This means that there are $3x$ people. Because no one's score was the same as another person's score, that means that there could only have been $1$ median score. This implies that $x$ is an odd number. $x$ cannot be less than $23$, because there wouldn't be a $64$th place if $x$ was. $x$ cannot be greater than $23$ either, because that would tie Andrea and Beth or Andrea's place would be worse than Beth's. Thus, the only possible answer is $\boxed{\mathbf{(B)}\ 23}$.

 



\textbf{Solution 2}

Let $a$ be Andrea's place. We know that she was the highest on her team, so $a < 37$. 

 Since $a$ is the median, there are $a-1$ to the left and right of the median, so the total number of people is $2a-1$ and the number of schools is $(2a-1)/3$. This implies that $2a-1 \equiv 0 \pmod{3} \implies a \equiv 2 \pmod{3}$. 

 Also, since $2a-1$ is the rank of the last-place person, and one of Andrea's teammates already got 64th place, $2a-1 > 64 \implies a \ge 33$.

 Putting it all together: $33 \le a < 37$ and $a \equiv 2 \pmod{3}$, so clearly $a = 35$, and the number of schools as we got before is $(2a-1)/3 = 69/3 = \boxed{\mathbf{(B)}\ 23}$.

 



\textbf{Solution 3}

Since all students have distinct scores, we can assume the scores are $1$ through $S$, where $S$ is the total number of students. From Solution 1, we know that $S$ is odd.

 Andrea’s score is the median, so$A = \frac{S+1}{2}$.Beth and Carla have scores$B = S - 37$, $\quad C = S - 64$.Since Andrea has the highest score among her teammates, we have$A > B \quad \text{and} \quad A > C$.Thus, $\frac{S+1}{2} > S - 37 \quad \text{and} \quad \frac{S+1}{2} > S - 64$.

 Solving these inequalities:$\frac{S+1}{2} > S - 37 \Rightarrow S < 75$,$\frac{S+1}{2} > S - 64 \Rightarrow S < 129$.Hence, $S < 75$.We also have : $B + C > A$.Substituting:$(S - 37) + (S - 64) > \frac{S+1}{2}$.This simplifies to$2S - 101 > \frac{S+1}{2}$$\Rightarrow 4S - 202 > S + 1$$\Rightarrow 3S > 203$$\Rightarrow S > \frac{203}{3}$.

 Since students are grouped into teams of 3, let$S = 3t$.Then$\frac{203}{3} < S < 75$.$\Rightarrow 22 < t < 25$Because t has to be an odd number in this range we get :$t = 23$.

 $\boxed{\mathbf{(B)\ 23}}$

 


	\item (\textit{AMC 10 10B 2010}) Positive integers $a$, $b$, and $c$ are randomly and independently selected with replacement from the set $\{1, 2, 3,\dots, 2010\}$. What is the probability that $abc + ab + a$ is divisible by $3$?


 $\textbf{(A)}\ \dfrac{1}{3} \qquad \textbf{(B)}\ \dfrac{29}{81} \qquad \textbf{(C)}\ \dfrac{31}{81} \qquad \textbf{(D)}\ \dfrac{11}{27} \qquad \textbf{(E)}\ \dfrac{13}{27}$


 
 \textbf{Answer: }E



\textbf{Solution 1}

We group this into groups of $3$, because $3|2010$. This means that every residue class mod 3 has an equal probability.

 If $3|a$, we are done. There is a probability of $\frac{1}{3}$ that that happens.

 Otherwise, we have $3|bc+b+1$, which means that $b(c+1) \equiv 2\pmod{3}$. So either 
 \[b \equiv 1 \pmod{3}, c \equiv 1 \pmod{3}\] or 
 \[b \equiv 2 \pmod {3}, c \equiv 0 \pmod 3\] which will lead to the property being true. There is a $\frac{1}{3}\cdot\frac{1}{3}=\frac{1}{9}$ chance for each bundle of cases to be true. Thus, the total for the cases is $\frac{2}{9}$. But we have to multiply by $\frac{2}{3}$ because this case ($3|bc+b+1$ ) only needs to happen when $a$ isn’t divisible by 3,  which has a $\frac{2}{3}$ chance. So the total is actually $\frac{4}{27}$.

 The grand total is 
 \[\frac{1}{3} + \frac{4}{27} = \boxed{\text{(E) }\frac{13}{27}.}\]

 



\textbf{Solution 2 (Minor change from Solution 1)}

Just like solution 1, we see that there is a $\frac{1}{3}$ chance of $3|a$ and $\frac{2}{9}$ chance of $3|1+b+bc$ 

 Now, we can just use PIE (Principles of Inclusion and Exclusion) to get our answer to be $\frac{1}{3}+\frac{2}{9}-\frac{1}{3}\cdot\frac{2}{9} = \boxed{\text{(E) } \frac{13}{27}}$

 -Conantwiz2023

 



\textbf{Solution 3 (Fancier version of Solution 1)}

As with solution one, we conclude that if $a\equiv0\mod 3$ then the requirements are satisfied. We then have:
 \[a(bc+c)+a\equiv0 \mod 3\]
 \[a(bc+c)\equiv-a \mod 3\]
 \[c(b+1)\equiv-1 \mod 3\]
 \[b+1\equiv \frac{-1}{c} \mod 3\]Which is true for one $b$ when $c\not\equiv 0 \mod 3$ because the integers$\mod 3$ form a field under multiplication and addition with absorbing element $0$.

 This gives us $P=\frac{1}{3}+\frac{4}{27}=\boxed{\text{(E) }\frac{13}{27}}$.

 


	\item (\textit{AMC 10 10B 2010}) A circle with center $O$ has area $156\pi$. Triangle $ABC$ is equilateral, $\overline{BC}$ is a chord on the circle, $OA = 4\sqrt{3}$, and point $O$ is outside $\triangle ABC$. What is the side length of $\triangle ABC$?


 $\textbf{(A)}\ 2\sqrt{3} \qquad \textbf{(B)}\ 6 \qquad \textbf{(C)}\ 4\sqrt{3} \qquad \textbf{(D)}\ 12 \qquad \textbf{(E)}\ 18$


 
 \textbf{Answer: }B



\textbf{Solution 1}

The formula for the area of a circle is $\pi r^2$ so the radius of this circle is $\sqrt{156}.$

 Because $OA=4\sqrt{3}=\sqrt{48} < \sqrt{156}, A$ must be in the interior of circle $O.$

 
                    \begin{figure}[H]
                        \centering
                        \includegraphics[width=0.5\textwidth]{Solution/AMC_10_10B_2010_19_Solution_Figure_1.png}
                    \end{figure}
                    

 Let $s$ be the side length of the triangle, the unknown value, and let $X$ be the point on $BC$ where $OX \perp BC.$ Since $\triangle ABC$ is equilateral, $BX=\frac{s}{2}$ and $AX=\frac{s\sqrt{3}}{2}.$ We are given $AO=4\sqrt{3}.$ Use the \href{https://artofproblemsolving.com/wiki/index.php?title=Pythagorean\_Theorem}{Pythagorean Theorem} and solve for $s.$

 
 \begin{align*} (\sqrt{156})^2 &= \left(\frac{s}{2}\right)^2 + \left( \frac{s\sqrt{3}}{2} + 4\sqrt{3} \right)^2\\ 156 &= \frac14s^2 + \frac34s^2 + 12s + 48\\ 0 &= s^2 + 12s - 108\\ 0 &= (s-6)(s+18)\\ s &= \boxed{\textbf{(B)}\ 6} \end{align*}


 



\textbf{Solution 2}

We can use the same diagram as Solution 1 and label the side length of $\triangle ABC$ as $s$. Using congruent triangles, namely the two triangles $\triangle BOA$ and $\triangle COA$, we get that $\angle BAO = \angle CAO \implies \angle BAO = \frac{360^\circ-60^\circ}{2} = 150^\circ$. From this, we can use the \href{https://artofproblemsolving.com/wiki/index.php?title=Law\_of\_Cosines}{Law of Cosines}, to get 
 \[s^2 + (4 \sqrt{3})^2 - 2 \times s \times 4 \sqrt{3} \times - \frac{\sqrt{3}}{2} = (2 \sqrt{39})^2\] Simplifying, we get 
 \[s^2 + 12s + 48 = 156 \implies s^2 + 12s - 108 = 0\] We can factor this to get 
 \[(x-6)(x+18)\] Lengths must be non-negative, so the answer is $\boxed{\textbf{(B)}\ 6}$~bryan gao

 


	\item (\textit{AMC 10 10B 2010}) Two circles lie outside regular hexagon $ABCDEF$. The first is tangent to $\overline{AB}$, and the second is tangent to $\overline{DE}$. Both are tangent to lines $BC$ and $FA$. What is the ratio of the area of the second circle to that of the first circle?


 $\textbf{(A)}\ 18 \qquad \textbf{(B)}\ 27 \qquad \textbf{(C)}\ 36 \qquad \textbf{(D)}\ 81 \qquad \textbf{(E)}\ 108$


 
 \textbf{Answer: }D



\textbf{Solution 1}

A good diagram is very helpful. 

  
                    \begin{figure}[H]
                        \centering
                        \includegraphics[width=0.5\textwidth]{Solution/AMC_10_10B_2010_20_Solution_Figure_1.gif}
                    \end{figure}
                     

 The first circle is in red, the second in blue.With this diagram, we can see that the first circle is inscribed in equilateral triangle $GBA$ while the second circle is inscribed in $GKJ$.From this, it's evident that the ratio of the blue area to the red area is equal to the \textbf{ratio of the areas $\triangle GKJ$ to $\triangle GBA$}

 Since the ratio of areas is equal to the square of the ratio of lengths, we know our final answer is $\left(\frac{GK}{GB}\right)^2$.From the diagram, we can see that this is $9^2=\boxed{\textbf{(D)}\ 81}$

 



\textbf{Solution 2}

As above, we note that the first circle is inscribed in an equilateral triangle of side length $1$ (if we assume, WLOG, that the regular hexagon has side length $1$). The inradius of an equilateral triangle with side length $1$ is equal to $\frac{\sqrt{3}}{6}$. Therefore, the area of the first circle is $(\frac{\sqrt{3}}{6})^2 \cdot \pi =\frac{\pi}{12}$.

 
Call the center of the second circle $O$. Now we drop a perpendicular from $O$ to circle $O$'s point of tangency with $GK$ and draw another line connecting $O$ to $G$. Note that because triangle $BGA$ is equilateral, $\angle BGA=60^{\circ}$. $OG$ bisects $\angle BGA$, so we have a $30-60-90$ triangle. 

 
Call the radius of circle $O$ $r$.$OG=2r= \text{height of equilateral triangle} + \text{height of regular hexagon} + r$

 The height of an equilateral triangle of side length $1$ is $\frac{\sqrt{3}}{2}$. The height of a regular hexagon of side length $1$ is $\sqrt{3}$. Therefore,$OG=\frac{\sqrt{3}}{2} + \sqrt{3} + r$.

 We can now set up the following equation:

 $\frac{\sqrt{3}}{2} + \sqrt{3} + r=2r$

 $\frac{\sqrt{3}}{2} + \sqrt{3}=r$

 $\frac{3\sqrt{3}}{2}=r$

 The area of Circle $O$ equals $\pi r^2=\frac{27}{4} \pi$

 Therefore, the ratio of the areas is $\frac{\frac{27}{4}}{\frac{1}{12}}=\boxed{\textbf{(D)}\ 81}$

 


	\item (\textit{AMC 10 10B 2010}) A palindrome between $1000$ and $10,000$ is chosen at random. What is the probability that it is divisible by $7$?


 $\textbf{(A)}\ \dfrac{1}{10} \qquad \textbf{(B)}\ \dfrac{1}{9} \qquad \textbf{(C)}\ \dfrac{1}{7} \qquad \textbf{(D)}\ \dfrac{1}{6} \qquad \textbf{(E)}\ \dfrac{1}{5}$


 
 \textbf{Answer: }E



\textbf{Solution 1}

View the palindrome as some number with form (decimal representation):$a_3 \cdot 10^3 + a_2 \cdot 10^2 + a_1 \cdot 10 + a_0$. But because the number is a palindrome, $a_3 = a_0, a_2 = a_1$. Recombining this yields $1001a_3 + 110a_2$. 1001 is divisible by 7, which means that as long as $a_2 = 0$, the palindrome will be divisible by 7. This yields 9 palindromes out of 90 ($9 \cdot 10$) possibilities for palindromes. However, if $a_2 = 7$, then this gives another case in which the palindrome is divisible by 7. This adds another 9 palindromes to the list, bringing our total to $18/90 = \boxed {\frac{1}{5} } = \boxed {E}$

 



\textbf{Solution (Divisibility Rules)}

We can notice the palindrome is of the form $\overline{abba}$. Then, by the divisibility rule of $7$, $7$ must divide 
 \[100a+11b-2a = 98a+11b.\] This nicely simplifies to the fact that $7 \mid 4b,$ so $b$ is clearly $0$ or $7$. This gives us $9 \cdot 2$ total choices for the palindrome divisible by $7$, divided by $9 \cdot 10$ total choices for $\overline{abba}$, giving us an answer of $\boxed{\text{(E)}} \ \dfrac{1}{5}$.

 \textbf{Addendum (Alternate)} $7\mid 1001a^3+110b^2$ and $1001 \equiv 0 \pmod 7$. Knowing that $a$ does not factor (pun intended) into the problem, note 110's prime factorization and $7\mid b$. There are only 10 possible digits for $b$, 0 through 9, but $7\mid b$ only holds if $b=0, 7$. This is 2 of the 10 digits, so $\frac{2}{10}=\boxed{\textbf{(E) }\frac{1}{5}}$

 



\textbf{Solution 4}

The palindrome is in the form of $\overline{abba}$. There are $9$ possible values of $a$ and 10 possible values for $b$. Thus, there are $90$ four-digit palindromes. The smallest palindrome is $1001$, which is a multiple of $7$. Note how each palindrome with the same thousands and ones digit increase by $110$ (i.e., $1111-1001=110$, $1221-1111=110$, etc). $110 \equiv 5$ (mod $7$). Since $5$ and $7$ is coprime, to get another palindrome that is a multiple of $7$, we need to increase by $770$. $1001+770=1771$. Next, when the thousands place is $2$, we start with $2002$, which is a multiple of $1001$, so $2002 \equiv 0$ (mod $7$). Therefore, $2772$ is also a satisfactory palindrome. Similarly, for every thousands place, there are $2$ satisfactory palindromes, which means there are a total of $2 \cdot 9 = 18$ satisfactory palindromes. Thus, the answer is $\frac{18}{90} = \frac{1}{5}$ or $\boxed{\text{(E)}\ \frac{1}{5}}$.

 


	\item (\textit{AMC 10 10B 2010}) Seven distinct pieces of candy are to be distributed among three bags. The red bag and the blue bag must each receive at least one piece of candy; the white bag may remain empty. How many arrangements are possible?


 $\textbf{(A)}\ 1930 \qquad \textbf{(B)}\ 1931 \qquad \textbf{(C)}\ 1932 \qquad \textbf{(D)}\ 1933 \qquad \textbf{(E)}\ 1934$


 
 \textbf{Answer: }C



\textbf{Solution 1}

We can count the total number of ways to distribute the candies (ignoring the restrictions), and then subtract the overcount to get the answer.

 Each candy has three choices; it can go in any of the three bags.

 Since there are seven candies, that makes the \textbf{total distribution} $3^7=2187$

 
To find the overcount, we calculate the number of invalid distributions: the red or blue bag is empty. 

 The number of distributions such that the red bag is empty is equal to $2^7$, since it's equivalent to distributing the $7$ candies into $2$ bags.

 We know that the number of distributions with the blue bag is empty will be the same number because of the symmetry, so it's also $2^7$. 

 The case where both the red and the blue bags are empty (all $7$ candies are in the white bag) are included in both of the above calculations, and this case has only $1$ distribution.

 \textbf{The total overcount is $2^7+2^7-1=2^8-1$}

 
The final answer will be $\text{total}-\text{overcount}=2187-(2^8-1) = 2187-256+1=1931+1=\boxed{\textbf{(C)}\ 1932}$

 You can do the last step slightly faster by noticing $3^7 \equiv 7 \pmod{10}$ and $2^8 \equiv 6 \pmod{10}$, so the end total will be $3^7 - 2^8 +1 \equiv 7 - 6 + 1 \equiv 2 \pmod{10}$. The only answer given that satisfies this is $\boxed{\textbf{(C)}\ 1932}$.

 



\textbf{Solution 2}

We can use to our advantage the answer choices $\text{AMC}$ has given us, and eliminate the obvious wrong answer choices.

 We can first figure out how many ways there are to take two candies from seven distinct candies to place them into the red/blue bags: $7\cdot 6=42$.

 Now we can look at the answer choices to find out which ones are divisible by $42$, since the total number of combinations must be $42$ multiplied by some other number.

 Since answers A, B, D, and E are not divisible by 3 (divisor of 42), the answer must be $\boxed{\textbf{(C)}\ 1932}$.

 



\textbf{Solution 3}

Let $r$ be the number of red bag candies. For ${1 \leq r \leq 6}$. 

 So the number of candies left for the blue bag and the red bag is $7-r$. Based on the problem, 1 candy must be fixed for the blue bag, which can be done $7-r$ ways. Now, before we continue, we need to realize that fixing a candy can lead to some over counting in cases where none in the white bag overlap. So we should try and find an alternative because we'll be over counting more than twice, and that will become extremely difficult to account for each case's over counting. 

 So, without fixing candies, we can put everything into options and work on the white bag, because once we figure out two bags, the remaining one is decided. 

 The options for the candies in the white bag are two: In the white bag or out of the white bag(by default, in the blue bag).

 Now, for choosing the red bag, we have ${7 \choose r}$.

 Then, we have for the white bag: $2^{7-r}$. Because we aren't fixing one for blue, the power is $7-r$ instead of $7-r-1$

 We have: $({7 \choose r}) \cdot (2^{7-r})$Adding them up for ${1 \leq r \leq 6}$, we get 2058.

 Now, for the invalid cases. Because we didn't fix candy for the blue bag, we need to subtract the cases where the blue bag remains empty. We can accomplish this pretty easily.

 When $r = 6$, how many ways can the remaining 1 candy $not$ be placed in the blue bag. This can be done ${7 \choose 1} = 7$ ways.

 When $r = 5$, how many ways can the remaining 2 candies $not$ be placed in the blue bag. This can be done ${7 \choose 2} = 21$ ways.

 When $r = 4$, how many ways can the remaining 3 candies $not$ be placed in the blue bag. This can be done ${7 \choose 3} = 35$ ways.

 And because ${7 \choose 4} = {7 \choose 3}, {7 \choose 5} = {7 \choose 2}, and {7 \choose 6} = {7 \choose 1}$, we just multiply by two. 

 Finally, we have: $2058 - 2 \cdot (7+21+35) = 2058 - 126 = \boxed{\textbf{(C)}\ 1932}$

 



\textbf{Solution 4 (easiest)}

If the candies were indistinct, we could simply use \href{https://artofproblemsolving.com/wiki/index.php?title=Ball-and-urn}{Stars and Bars}. Because the candies are distinct, you will have to multiply by $7!$ to get the correct answer. Knowing that $7!$ is a multiple of $7$, the only answer that is a multiple of 7 is $\boxed{\textbf{(C)}\ 1932}$

 


	\item (\textit{AMC 10 10B 2010}) The entries in a $3 \times 3$ array include all the digits from $1$ through $9$, arranged so that the entries in every row and column are in increasing order. How many such arrays are there?


 $\textbf{(A)}\ 18 \qquad \textbf{(B)}\ 24 \qquad \textbf{(C)}\ 36 \qquad \textbf{(D)}\ 42 \qquad \textbf{(E)}\ 60$


 
 \textbf{Answer: }D



\textbf{Solution 1}

Observe that all tables must have 1s and 9s in the corners, 8s and 2s next to those corner squares, and 4-6 in the middle square. Also note that for each table, there exists a valid table diagonally symmetrical across the diagonal extending from the top left to the bottom right. 

 


 
                \begin{itemize}
                    \item \textbf{Case 1: Center 4}
                \end{itemize}


                
 \[\begin{tabular}{|c|c|c|} \hline 1&2&\\ \hline 3&4&8\\ \hline &&9\\ \hline \end{tabular} \;\;\; \begin{tabular}{|c|c|c|} \hline 1&2&\\ \hline 3&4&\\ \hline &8&9\\ \hline \end{tabular}\]

 3 necessarily must be placed as above. Any number could fill the isolated square, but the other 2 are then invariant. So, there are 3 cases each and 6 overall cases. Given diagonal symmetry, alternate 2 and 8 placements yield symmetrical cases. $2*6=12$

 
                \begin{itemize}
                    \item \textbf{Case 2: Center 5}
                \end{itemize}


                
 \[\begin{tabular}{|c|c|c|} \hline 1&2&3\\ \hline 4&5&\\ \hline &8&9\\ \hline \end{tabular} \;\;\; \begin{tabular}{|c|c|c|} \hline 1&2&\\ \hline 3&5&\\ \hline &8&9\\ \hline \end{tabular} \;\;\; \begin{tabular}{|c|c|c|} \hline 1&2&\\ \hline 3&5&8\\ \hline &&9\\ \hline \end{tabular} \;\;\; \begin{tabular}{|c|c|c|} \hline 1&2&3\\ \hline 4&5&8\\ \hline &&9\\ \hline \end{tabular}\]

 Here, no 3s or 7s are assured, but this is only a teensy bit trickier and messier. WLOG, casework with 3 instead of 7 as above. Remembering that $4<5$, logically see that the numbers of cases are then 2,3,3,1 respectively. By symmetry, $2*9=18$

 
                \begin{itemize}
                    \item \textbf{Case 3: Center 6}
                \end{itemize}


                By inspection, realize that this is symmetrical to case 1 except that the 7s instead of the 3s are assured.$2*6=12$

 
 \[12+18+12=\boxed{\textbf{(D) }42}\]

 



\textbf{Solution 2}

This solution is trivial by the hook length theorem. The hooks look like this:

 $\begin{tabular}{|c|c|c|} \hline 5 & 4 & 3 \\ \hline 4 & 3 & 2\\ \hline 3 & 2 & 1\\ \hline \end{tabular}$

 So, the answer is $\frac{9!}{5 \cdot 4 \cdot 3 \cdot 4 \cdot 3 \cdot 2 \cdot 3 \cdot 2 \cdot 1}$ = $\boxed{\text{(D) }42}$

 


	\item (\textit{AMC 10 10B 2010}) A high school basketball game between the Raiders and Wildcats was tied at the end of the first quarter. The number of points scored by the Raiders in each of the four quarters formed an increasing geometric sequence, and the number of points scored by the Wildcats in each of the four quarters formed an increasing arithmetic sequence. At the end of the fourth quarter, the Raiders had won by one point. Neither team scored more than $100$ points. What was the total number of points scored by the two teams in the first half?


 $\textbf{(A)}\ 30 \qquad \textbf{(B)}\ 31 \qquad \textbf{(C)}\ 32 \qquad \textbf{(D)}\ 33 \qquad \textbf{(E)}\ 34$


 
 \textbf{Answer: }E



\textbf{Solution 1}

Let $a,ar,ar^{2},ar^{3}$ be the quarterly scores for the Raiders. We know $r > 1$ because the sequence is said to be increasing. We also know that each of $a, ar, ar^2, ar^3$ is an integer. We start by showing that \textbf{$r$ must also be an integer.}

 Suppose not, and say $r = m/n$ where $m>n>1$, and $\gcd(m,n)=1$. Then $n, n^2, n^3$ must all divide $a$ so $a=n^3k$ for some integer $k$. Then $S_R = n^3k + n^2mk + nm^2k + m^3k < 100$ and we see that even if $k=1$ and $n=2$, we get $m < 4$, which means that the only option for $r$ is $r=3/2$. A quick check shows that even this doesn't work. Thus $r$ must be an integer.

 Let $a, a+d, a+2d, a+3d$ be the quarterly scores for the Wildcats. Let $S_W = a+(a+d) + (a+2d)+(a+3d) = 4a+6d$. Let $S_R = a+ar+ar^2+ar^3 = a(1+r)(1+r^2)$. Then $S_R<100$ implies that $r<5$, so $r\in \{2, 3, 4\}$. The Raiders win by one point, so
 \[a(1+r)(1+r^2) = 4a+6d+1.\]

 
                \begin{itemize}
                    \item If $r=4$ we get $85a = 4a+6d+1$ which means $3(27a-2d) = 1$, which is not possible with the given conditions.\item If $r=3$ we get $40a = 4a+6d+1$ which means $6(6a-d) = 1$, which is also not possible with the given conditions.\item If $r=2$ we get $15a = 4a+6d+1$ which means $11a-6d = 1$. Reducing modulo 6 we get $a \equiv 5\pmod{6}$. Since $15a<100$ we get $a<7$. Thus $a=5$. It then follows that $d=9$.
                \end{itemize}


                Then the quarterly scores for the Raiders are $5, 10, 20, 40$, and those for the Wildcats are $5, 14, 23, 32$. Also $S_R = 75 = S_W + 1$. The total number of points scored by the two teams in the first half is $5+10+5+14=\boxed{\textbf{(E)}\ 34}$.

 Note if you don't realize while taking the test that $r$ might not be an integer: since an answer is achieved through casework on the integer value of $r$ and since there is only one right answer, the proof of $r$ being an integer can be skipped on the test (it takes up time).

 



\textbf{Solution 2}

Let $a,ar,ar^{2},ar^{3}$ be the quarterly scores for the Raiders. We know that the Raiders and Wildcats both scored the same number of points in the first quarter so let $a,a+d,a+2d,a+3d$ be the quarterly scores for the Wildcats. The sum of the Raiders scores is $a(1+r+r^{2}+r^{3})$ and the sum of the Wildcats scores is $4a+6d$. Now we can narrow our search for the values of $a,d$, and $r$. Because points are always measured in positive integers, we can conclude that $a$ and $d$ are positive integers. We can also conclude that $r$ is a positive integer by writing down the equation:

 
 \[a(1+r+r^{2}+r^{3})=4a+6d+1\]

 Now we can start trying out some values of $r$. We try $r=2$, which gives

 
 \[15a=4a+6d+1\]

 
 \[11a=6d+1\]

 We need the smallest multiple of $11$ (to satisfy the <100 condition) that is $\equiv 1 \pmod{6}$. We see that this is $55$, and therefore $a=5$ and $d=9$.

 So the Raiders' first two scores were $5$ and $10$ and the Wildcats' first two scores were $5$ and $14$. 

 
 \[5+10+5+14=34 \longrightarrow \boxed{\textbf{(E)}}\]

 



\textbf{Solution 3 (Quick Solve?)}

It should become apparent that the geometric ratio is less than or equal to $2$ (try the first quarter $1$, ratio $3$). Obviously, the ratio is more than $1$, so we try the ratio of $2$

 $1, 2, 4, 8$ - The Raiders 

 $1, ?, ?, ?$ - The Wildcats

 To find the $?$s, we need the sum of the raiders, which is $15$. $15-1=14$, and $\frac{14}2=7$. So the last term is $7-1=6$. But, now we can't have integrated 2nd and 3rd quarter scores, so we stop.

 $2, 4, 8, 16$ - The Raiders, this also fails (try it out)

 The same goes for $3, 6, 12, 24$, and $4, 8, 16, 32$. With $5, 10, 20, 40$ we get:

 $5, 10, 20, 40$ - The Raiders

 $5, ?, ?, ?$ - The wildcats

 $5+10+20+40=75$, $75-1=74$, $\frac{74}2=37$, $37-5=32$

 $5, ?, ?, 32$ - The Wildcats.

 Now, we can finish the arithmetic progression:

 $5, 14, 23, 32$ - The Wildcats

 Thus our answer is $5+10+5+14=34$

 



\textbf{Solution 4 (solution 2 variant)}

Raiders: $a, am, am^2, am^3$ 

 Wildcats: $a, a + n, a + 2n, a + 3n,$ 

 where $a$ represents both teams' score in the first quarter. Let the Raiders finish with a score of $r$. We know $4a + 6n + 1 = a(1 + m + m^2 + m^3) = r.$ Since $r \le 100$, $1 + m + m^2 + m^3$ can either be $4, 15, 40,$ or $85$. Solve for $a$ in both equations in terms of $r$ and set them equal to each other:

 $\frac{r - (6n + 1)}{4} = \frac{r}{1 + m + m^2 + m^3}.$ We know $1 + m + m^2 + m^3$ can't be $4$ because then $n$ would be a fraction. So let's try $15$:

 $\frac{r - (6n + 1)}{4} = \frac{r}{15} \implies 11r = 90n + 15.$ We know $r$ has to be a multiple of 15 for $a$ to be an integer. After some guess and check with the first few multiples of 15, we find that if $r = 75, n = 9.$ Then $a = 5$. Since this whole case is based on $1 + m + m^2 + m^3 = 15, m = 2$. So our answer is $3a + am + n = \boxed{34}.$

 \textbf{ᴠɪᴅᴇᴏ ѕᴏʟᴜᴛɪᴏɴ ʙʏ ᴘɪ ᴀᴄᴀᴅᴇᴍʏ} 


	\item (\textit{AMC 10 10B 2010}) Let $a > 0$, and let $P(x)$ be a polynomial with integer coefficients such that


 
$P(1) = P(3) = P(5) = P(7) = a$, and

$P(2) = P(4) = P(6) = P(8) = -a$.


 


 What is the smallest possible value of $a$?


 $\textbf{(A)}\ 105 \qquad \textbf{(B)}\ 315 \qquad \textbf{(C)}\ 945 \qquad \textbf{(D)}\ 7! \qquad \textbf{(E)}\ 8!$


 
 \textbf{Answer: }B



\textbf{Solution 1}

We observe that because $P(1) = P(3) = P(5) = P(7) = a$, if we define a new polynomial $R(x)$ such that $R(x) = P(x) - a$, $R(x)$ has roots when $P(x) = a$; namely, when $x=1,3,5,7$.

 Thus since $R(x)$ has roots when $x=1,3,5,7$, we can factor the product $(x-1)(x-3)(x-5)(x-7)$ out of $R(x)$ to obtain a new polynomial $Q(x)$ such that $(x-1)(x-3)(x-5)(x-7)(Q(x)) = R(x) = P(x) - a$.

 Then, noting that $P(2)-a=-a-a=-2a$ etc., and plugging in values of $2,4,6,8,$ we get 

 
 \[P(2)-a=(2-1)(2-3)(2-5)(2-7)Q(2) = -15Q(2) = -2a\]
 \[P(4)-a=(4-1)(4-3)(4-5)(4-7)Q(4) = 9Q(4) = -2a\]
 \[P(6)-a=(6-1)(6-3)(6-5)(6-7)Q(6) = -15Q(6) = -2a\]
 \[P(8)-a=(8-1)(8-3)(8-5)(8-7)Q(8) = 105Q(8) = -2a\]

 $-2a=-15Q(2)=9Q(4)=-15Q(6)=105Q(8).$Thus, the least value of $a$ must be the $\text{lcm}(15,9,15,105)$.Solving, we receive $315$, so our answer is $\boxed{\textbf{(B)}\ 315}$.

 To complete the solution, we can let $a = 315$, and then try to find $Q(x)$. We know from the above calculation that $Q(2)=42, Q(4)=-70, Q(6)=42$, and $Q(8)=-6$. Then we can let $Q(x) = T(x)(x-2)(x-6)+42$, getting $T(4)=28, T(8)=-4$. Let $T(x)=L(x)(x-8)-4$, then $L(4)=-8$. Therefore, it is possible to choose $T(x) = -8(x-8)-4 = -8x + 60$, so the goal is accomplished. As a reference, the polynomial we get is

 
 \[P(x) = (x-1)(x-3)(x-5)(x-7)((-8x + 60)(x-2)(x-6)+42) + 315\]
 \[= -8 x^7+252 x^6-3248 x^5+22050 x^4-84392 x^3+179928 x^2-194592 x+80325\]

 



\textbf{Solution 2}

The evenly-spaced data suggests using \href{https://artofproblemsolving.com/wiki/index.php?title=Discrete\_derivative}{discrete derivatives} to tackle this problem.  First, note that any polynomial of degree $n$

 $P(x) = a_0 + a_1 x + a_2 x^2 + \ldots + a_n x^n$

 can also be written as

 $P(x) = b_0 + b_1 (x-1) + b_2 (x-1)(x-2) + \ldots + b_n (x-1)(x-1) \cdots (x-n)$.

 Moreover, the coefficients $a_i$ are integers for $i=1, 2, \ldots n$ iff the coefficients $b_i$ are integers for $i=1, 2, \ldots n$.  This latter form is convenient for calculating discrete derivatives of $P(x)$.

 The discrete derivative of a function $f(x)$ is the related function $\Delta f(x)$ defined as

 $\delta f(x) = f(x+1) - f(x)$.

 With this definition, it's easy to see that for any positive integer $k$ we have

 $\Delta [(x-1)(x-2)\cdots(x-k)] = k(x-1)(x-2)\cdots(x-[k-1])$.

 This in turn allows us to use successive discrete derivatives evaluated at $x=1$ to calculate all of the coefficients $b_i$ using

 $P(1) = b_0$, $\Delta P(1) = b_1$, $\Delta^2 P(1) = 2 b_2$, $\ldots$, $\Delta^7 P(1) = 7! b_7$.

 We can also calculate the following table of discrete derivatives based on the data points given in the problem statement:

 <table cellpadding="3" frame="box" rules="all"><tbody><tr><th></th><th colspan="8">$x$</th></tr><tr><td align="right"></td><td align="center">$1$</td><td align="center">$2$</td><td align="center">$3$</td><td align="center">$4$</td><td align="center">$5$</td><td align="center">$6$</td><td align="center">$7$</td><td align="center">$8$</td></tr><tr><td align="right">$P(x)$</td><td align="right">$a$</td><td align="right">$-a$</td><td align="right">$a$</td><td align="right">$-a$</td><td align="right">$a$</td><td align="right">$-a$</td><td align="right">$a$</td><td align="right">$-a$</td></tr><tr><td align="right">$\Delta P(x)$</td><td align="right">$-2a$</td><td align="right">$2a$</td><td align="right">$-2a$</td><td align="right">$2a$</td><td align="right">$-2a$</td><td align="right">$2a$</td><td align="right">$-2a$</td><td></td></tr><tr><td align="right">$\Delta^2 P(x)$</td><td align="right">$4a$</td><td align="right">$-4a$</td><td align="right">$4a$</td><td align="right">$-4a$</td><td align="right">$4a$</td><td align="right">$-4a$</td><td></td><td></td></tr><tr><td align="center" colspan="9">$\vdots$</td></tr><tr><td align="right">$\Delta^7 P(x)$</td><td align="right">$-2^7 a$</td><td></td><td></td><td></td><td></td><td></td><td></td><td></td></tr></tbody></table>

 Thus we can read down the column for $x=1$ to find that $k! b_k = (-2)^k a$ for $k = 0, 1, \ldots, 7$.  Interestingly, even if we choose $P(x)$ to have degree greater than $7$, the $8$ coefficients of lowest order always satisfy these conditions.  Moreover, it's straightforward to show that the $P(x)$ of degree $7$ with $b_k$ satisfying these conditions will fit the data given in the problem statement.  Thus, to find minimal necessary and sufficient conditions on the value of $a$, we need only consider these $8$ equations.  As a result, $P(x)$ with integer coefficients fitting the given data exists iff $k!$ divides $2^k a$ for $k = 0, 1, \ldots, 7$.  In other words, it's necessary and sufficient that

 
 \[0! | a\],

 
 \[1! | 2a\],

 
 \[2! | 2^2 a\],

 
 \[3! | 2^3 a\],

 
 \[4! | 2^4 a\],

 
 \[5! | 2^5 a\],

 
 \[6! | 2^6 a, \text{and}\]

 
 \[7! | 2^7 a.\]

 The last condition holds if $7 \cdot 3 \cdot 5 \cdot 3 = 315$ divides evenly into $a$.  Since such $a$ will also satisfy the first $7$ conditions, our answer is $\boxed{\textbf{(B)}\ 315}$.

 



\textbf{Solution 3 (2 liner)}

We have $P(x)-a=(x-1)(x-3)(x-5)(x-7)Q(x).$ Plugging in $x=2,4,6,$ and $8$, we get
 \[Q(2)=\dfrac{2}{15}a\]
 \[Q(4)=-\dfrac{2}{9}a\]
 \[Q(6)=-\dfrac{2}{15}a\]
 \[Q(8)=\dfrac{2}{105}a.\]Thus, $a=\text{lcm}(9,15,105)=\boxed{\textbf{(B)}\ 315}.$

 -jb2015007

 


\end{enumerate}